\chapter{Construire un lecteur, de A à Z}

Un lecteur est la première application où vous ne \emph{décidez} de rien : vous
ouvrez un fichier, vous le transformez, vous l'affichez. Tout le travail est dans
la transformation --- et c'est justement là que se cachent les erreurs qui
donnent une image inclinée, une forme d'onde plate ou une vidéo qui saccade.

Nous construisons \textbf{trois} lecteurs --- images, son, vidéo --- avec le même
squelette. Neuf étapes, et vous les avez tous les trois.

\begin{center}
\begin{tabular}{@{}llr@{}}
\toprule
& \textbf{Ce qu'il lit} & \textbf{Lignes} \\
\midrule
\texttt{NkAudioPlayer} & un son, et sa forme d'onde & 263 \\
\texttt{NkVideoPlayer} & une vidéo, et son audio & 684 \\
\texttt{NkImageDemo} & une image, et \emph{ce qu'il ne faut pas faire} & 2\,369 \\
\bottomrule
\end{tabular}
\end{center}

\begin{nkretenir}
\textbf{Les étapes 1 et 2 du chapitre 21 ne changent pas} : le dossier, le
\texttt{.jenga}, le \texttt{main.cpp} de douze lignes. Ce qui change commence à
la troisième --- et ce chapitre ne réécrit que ce qui change.

Aux dépendances du \texttt{.jenga}, ajoutez selon le lecteur :
\texttt{NKImage} (déjà là), \texttt{NKAudio} (déjà là), et
\textbf{\texttt{NKMedia}} pour la vidéo.
\end{nkretenir}

\section{Étape 1 --- ouvrir un fichier, et le dire quand ça rate}

\begin{nkcode}{Le squelette d'ouverture, commun aux trois}
bool OnGuiInit() override {
    const NkString chemin = mArguments.Size() > 1 ? mArguments[1] : NkString();
    if (chemin.Empty()) {
        mMessage = "Usage : NkLecteur <fichier>";
        return true;                       // on demarre quand meme, et on le DIT
    }
    if (!Ouvrir(chemin)) {
        mMessage = NkFormat("Impossible de lire : {0}", chemin);
        return true;                       // fenetre ouverte, message lisible
    }
    return true;
}
\end{nkcode}

\begin{nkretenir}
\textbf{On démarre même quand le fichier manque.} Rendre \texttt{false} fermerait
l'application avant qu'un message soit lisible --- l'utilisateur voit une fenêtre
qui clignote et disparaît, et il n'a aucun moyen d'apprendre pourquoi.

\textbf{Un repli se crie, mais ne verrouille pas l'utilisateur dehors.}
\end{nkretenir}

\begin{nkpiege}
\textbf{Ne déduisez pas le format de l'extension.} Un \texttt{.jpg} peut être un
PNG renommé, et un \texttt{.avi} peut contenir n'importe quoi.

Pour la vidéo, NKMedia offre \texttt{NkMediaProbe} : on \emph{sonde} le fichier
avant de l'ouvrir. Pour l'image, le décodeur lit la signature des premiers
octets. Dans les deux cas, \textbf{l'en-tête est la vérité, l'extension n'est
qu'un indice.}
\end{nkpiege}

\section{Étape 2 --- le son : réduire un million d'échantillons}

Un fichier de trois minutes en 44\,100 Hz contient environ huit millions
d'échantillons. Votre fenêtre en fait mille de large. Il faut donc en montrer
huit mille par colonne --- et le \emph{comment} n'est pas un détail.

\begin{nkcode}{Applications/NkAudioPlayer/src/main.cpp}
static void ComputePeaks(const audio::AudioSample &s, int32 cols,
                         NkVector<WavePeak> &out) {
    out.Clear();
    if (!s.data || s.frameCount <= 0 || cols <= 0) { return; }
    out.Resize((uint64)cols);
    const int32 ch = s.channels > 0 ? s.channels : 1;

    for (int32 c = 0; c < cols; ++c) {
        const uint64 f0 = (uint64)((double)c / cols * (double)s.frameCount);
        uint64 f1 = (uint64)((double)(c + 1) / cols * (double)s.frameCount);
        if (f1 <= f0) { f1 = f0 + 1; }
        if (f1 > (uint64)s.frameCount) { f1 = (uint64)s.frameCount; }

        float32 lo = 1.0f, hi = -1.0f;
        for (uint64 f = f0; f < f1; ++f) {
            float32 m = 0.0f;
            for (int32 k = 0; k < ch; ++k) { m += s.data[f * (uint64)ch + (uint64)k]; }
            m /= (float32)ch;
            if (m < lo) { lo = m; }
            if (m > hi) { hi = m; }
        }
        if (hi < lo) { lo = 0.0f; hi = 0.0f; }
        out[(uint64)c].lo = lo;
        out[(uint64)c].hi = hi;
    }
}
\end{nkcode}

\begin{nkretenir}
\textbf{Chaque colonne garde le minimum et le maximum de sa tranche}, puis se
dessine comme un trait vertical de l'un à l'autre.

\textbf{Pourquoi pas la moyenne :} une onde sonore oscille autour de zéro. La
moyenne d'un passage fort et celle d'un silence sont toutes deux proches de
zéro --- vous obtiendriez une ligne plate, techniquement juste et parfaitement
inutile.

Ce qui porte l'information, c'est l'\textbf{amplitude}, donc l'écart. C'est un
cas d'école du piège général : \emph{un agrégat peut être exact et ne rien
décrire.}
\end{nkretenir}

\begin{nkpiege}
Trois détails de cette fonction, et chacun est un défaut évité :

\textbf{\texttt{lo = 1.0f, hi = -1.0f}} --- à l'envers des bornes attendues.
C'est la seule initialisation correcte d'une recherche de minimum et de maximum :
la première valeur rencontrée les remplace tous les deux. Initialiser à zéro
serait faux --- un signal entièrement négatif rendrait \texttt{hi = 0}, un
maximum qui n'existe pas dans les données.

\textbf{\texttt{if (f1 <= f0) f1 = f0 + 1;}} --- sur un fichier plus court que la
fenêtre, deux colonnes tomberaient sur la même tranche vide.

\textbf{\texttt{if (hi < lo) \{ lo = 0; hi = 0; \}}} --- la tranche vide. Sans
elle, une colonne sans échantillon garderait $1$ et $-1$, et dessinerait un
trait pleine hauteur à partir de rien.
\end{nkpiege}

\section{Étape 3 --- dessiner, et comprendre ce qu'on dessine}

\begin{nkcode}{Applications/NkAudioPlayer/src/main.cpp}
static void PushQuad(NkVector<NkVertex> &v, float32 x0, float32 y0,
                     float32 x1, float32 y1, const NkColor2D &col) {
    NkVertex a{x0, y0, 0, 0, col.r, col.g, col.b, col.a};
    NkVertex b{x1, y0, 0, 0, col.r, col.g, col.b, col.a};
    NkVertex c{x1, y1, 0, 0, col.r, col.g, col.b, col.a};
    NkVertex d{x0, y1, 0, 0, col.r, col.g, col.b, col.a};
    v.PushBack(a); v.PushBack(b); v.PushBack(c);
    v.PushBack(a); v.PushBack(c); v.PushBack(d);
}
\end{nkcode}

\begin{nkretenir}
\textbf{Quatre coins, six sommets.} Le GPU ne connaît que des triangles : un
rectangle est le triangle $abc$ plus le triangle $acd$. Les sommets $a$ et $c$
sont donc envoyés \emph{deux fois} --- c'est normal, c'est le prix d'une liste de
triangles.

C'est exactement ce que fait \texttt{AddRectFilled} pour vous. Le voir écrit une
fois explique tout le reste : pourquoi un dessin coûte en \emph{nombre de
sommets}, et pourquoi regrouper les rectangles dans un même tampon est ce qui
rend un rendu rapide.
\end{nkretenir}

\begin{nkretenir}
\textbf{Mille colonnes, c'est mille rectangles --- soit un seul appel de dessin}
si vous les poussez tous dans le même tampon avant de vider. Une forme d'onde
n'est donc pas coûteuse, contrairement à ce qu'on craint.
\end{nkretenir}

\section{Étape 4 --- la vidéo : la chaîne complète}

Six lignes, et elles sont dans l'en-tête du lecteur du dépôt :

\begin{nkcode}{Applications/NkVideoPlayer/src/main.cpp}
//   1. NkWindow          -> fenetre native (Win/Linux/macOS/...)
//   2. NkRenderWindow    -> cible de rendu 2D NKCanvas
//   3. NkVideoReader     -> decode chaque image en RGBA
//   4. NkTexture         -> recoit les pixels RGBA de l'image courante (Update)
//   5. NkSprite          -> affiche la texture, mise a l'echelle
//   6. boucle cadencee au fps de la video (Clear -> Draw -> Display)
\end{nkcode}

\begin{nkretenir}
\textbf{Il n'y a aucune API « vidéo » dans le rendu.} Une vidéo, pour l'écran,
c'est une texture qu'on remplace à chaque image.

Le décodage --- H.264, MJPEG, conteneurs --- est entièrement dans NKMedia, et le
rendu n'en sait rien. C'est la séparation la plus profitable de tout ce cours :
\emph{le module qui produit des pixels ignore le module qui les montre}.
\end{nkretenir}

\begin{nkcode}{Applications/NkVideoPlayer/src/main.cpp}
media::NkVideoReader reader;
// ... ouverture ...
const media::NkVideoReaderInfo &vin = reader.Info();   // dimensions, fps, duree

// puis, une fois par image :
frameTex.Update(fr.rgba.Data(), (uint32)fr.width, (uint32)fr.height, 0, 0);
\end{nkcode}

\begin{nkretenir}
\textbf{Une seule ligne de téléversement, appelée une fois par image.} Tout le
reste du fichier --- 684 lignes --- est de la cadence, de la navigation et de
l'étalonnage de couleur.
\end{nkretenir}

\section{Étape 5 --- la cadence, et pourquoi elle n'est pas un \texttt{Sleep}}

\begin{nkcode}{Le principe}
mHorloge += deltaTime;
const float32 periode = 1.f / vin.fps;
while (mHorloge >= periode && reader.ReadFrame(fr)) {
    mHorloge -= periode;
    frameTex.Update(fr.rgba.Data(), fr.width, fr.height, 0, 0);
}
\end{nkcode}

\begin{nkretenir}
\textbf{Un \texttt{while}, pas un \texttt{if}.} Si une image a mis trop longtemps
à arriver --- une image-clé coûteuse, un coup de charge système --- il faut en
rattraper plusieurs d'un coup, sinon la vidéo prend un retard qu'elle ne
rattrapera jamais.

Et \texttt{mHorloge -= periode} plutôt que \texttt{mHorloge = 0} : remettre à
zéro perdrait le reste et ferait dériver la lecture de quelques pour cent.
\end{nkretenir}

\begin{nkpiege}
\textbf{N'attendez jamais avec un \texttt{Sleep} dans la boucle.} Le lecteur
cesserait de répondre : pas de redimensionnement, pas de clic, et sur le Web
l'onglet gèle.

La cadence se tient en \emph{comptant le temps écoulé}, pas en le dépensant.
\end{nkpiege}

\section{Étape 6 --- traiter les pixels, et à quel endroit}

Le lecteur du dépôt applique un étalonnage --- noir et blanc, sépia, table
\texttt{.cube} --- et l'endroit où il le fait est un enseignement à lui seul :

\begin{nkcode}{Applications/NkVideoPlayer/src/main.cpp (commentaire d'origine)}
// cablage est le meme : transformer les pixels juste avant frameTex.Update
\end{nkcode}

\begin{nkretenir}
\textbf{Un traitement d'image se branche entre le décodeur et la texture.} Ni
avant --- le décodeur ne doit rien savoir de vos couleurs --- ni après : une fois
téléversés, les pixels sont sur le GPU et ne vous appartiennent plus.

Cette position est un \emph{point d'extension} : ajouter un filtre ne demande de
toucher à rien d'autre. Écrivez-le en commentaire à cet endroit, c'est ce qui
guidera le prochain lecteur du fichier.
\end{nkretenir}

\section{Étape 7 --- l'image : le pas de ligne}

\begin{nkretenir}
\textbf{Le pas de ligne (\emph{stride}) n'est pas la largeur en octets.} Un
décodeur peut aligner chaque ligne sur quatre octets : une image de 641 pixels en
RGB occupe $641 \times 3 = 1923$ octets utiles, mais $1924$ par ligne.

Copier en supposant « largeur $\times$ octets » décale chaque ligne d'un octet de
plus que la précédente : \textbf{l'image sort inclinée}, en escalier, avec les
couleurs mélangées.
\end{nkretenir}

\begin{nkpiege}
\textbf{Le défaut est intermittent, et c'est ce qui le rend cher.} Il ne se voit
que si la largeur n'est pas un multiple de l'alignement --- donc jamais sur une
image de 640 ou 1024 pixels, et toujours sur celle de l'utilisateur.

\emph{Corriger en redimensionnant l'image à une largeur multiple de quatre
« marche » et est le pire des correctifs} : il masque la cause et la laisse en
place pour le prochain format.
\end{nkpiege}

\section{Étape 8 --- rendre les pixels, et les quatre cas}

\begin{nkretenir}
\textbf{La question n'est jamais « comment libère-t-on ceci ? » mais « qui l'a
alloué ? ».}

Quatre cas, et ils ne se devinent pas : une \texttt{NkImage} sur la pile se
détruit seule ; une allouée par \texttt{Alloc()} se rend par \texttt{Free()}
\emph{sans} \texttt{delete} ; un tampon rendu par un encodeur vient de
\texttt{NkAlloc} et se rend par \texttt{NkFree} ; et \texttt{Free()} appelé sur
un objet de la pile exécute une libération \emph{sur une adresse de pile}.

Ce dernier a réellement fait planter une démo du dépôt.
\end{nkretenir}

\section{Étape 9 --- la barre de transport}

\begin{nkcode}{Ce qu'elle doit porter, et rien de plus}
// [<<]  [ lecture / pause ]  [>>]   ####-------------  01:23 / 04:56
\end{nkcode}

\begin{nkretenir}
\textbf{La position affichée vient du lecteur, jamais d'un compteur à vous.} Un
compteur local dérive dès la première image sautée, et l'écart grandit sans
jamais se corriger.

Après un déplacement dans le fichier, c'est encore le lecteur qui dit où il est
réellement arrivé --- il se cale sur l'image-clé précédente, qui n'est pas
l'endroit demandé.
\end{nkretenir}

\section{Ce que \texttt{NkImageDemo} enseigne malgré lui}

\texttt{NkImageDemo} fait 2\,369 lignes pour afficher une image --- neuf fois le
lecteur audio. La raison est visible dans ses fichiers :

\begin{center}
\begin{tabular}{@{}ll@{}}
\toprule
\textbf{Fichier} & \textbf{Ce que c'est} \\
\midrule
\texttt{Render/GLContext} & un contexte OpenGL, à la main \\
\texttt{Render/GLRenderer2D} & un renderer 2D, à la main \\
\texttt{Render/Texture2D} & une texture, à la main \\
\texttt{Render/FontAtlas} & un atlas de police, à la main \\
\texttt{Render/SafeArea} & la zone sûre, à la main \\
\texttt{ViewerApp} & le visualiseur proprement dit \\
\bottomrule
\end{tabular}
\end{center}

\begin{nkretenir}
\textbf{Cinq de ses six fichiers réimplémentent ce que NKCanvas porte déjà.}

C'est le doublon dans sa forme la plus coûteuse : le code marche, personne n'a
tort, et pourtant chaque correction apportée à \texttt{NkRenderWindow} ne
parviendra jamais à \texttt{GLRenderer2D}.

\textbf{Avant d'écrire un mécanisme, cherchez qui le porte déjà --- et commencez
par la couche du dessous.} Un \texttt{grep} lancé dans son propre dossier ne
trouve jamais ce que la bibliothèque porte.
\end{nkretenir}

\begin{nkpiege}
\textbf{Les trois lecteurs du dépôt sont antérieurs à \texttt{NkCanvasApp}.} Ils
ouvrent la fenêtre eux-mêmes, dans \texttt{nkmain}, et tiennent leur propre
boucle.

Ce n'est pas une raison de les ignorer --- ils montrent exactement ce que la
coquille vous épargne --- mais c'en est une de ne pas les recopier tels quels.
Construisez les vôtres sur \texttt{NkCanvasGuiApp}, comme aux étapes 1 et 2 du
chapitre 21.
\end{nkpiege}

\begin{nkbilan}
Un lecteur ne décide de rien : tout son travail est une transformation. On ouvre
en sondant le contenu, pas l'extension, et l'on démarre même quand ça rate ---
avec un message lisible. Réduire un signal pour l'afficher demande le
\emph{minimum et le maximum}, jamais la moyenne. Un rectangle plein est deux
triangles et six sommets. Une vidéo, pour le rendu, n'est qu'une texture qu'on
remplace, cadencée par un \texttt{while} qui rattrape et jamais par un
\texttt{Sleep} ; le traitement des pixels se branche entre le décodeur et la
texture. Le pas de ligne n'est pas la largeur. Et \texttt{NkImageDemo} montre ce
que coûte un doublon : 2\,369 lignes dont cinq sixièmes existaient déjà un étage
plus bas.
\end{nkbilan}

\begin{nkquestions}
\begin{enumerate}
\item Pourquoi démarre-t-on l'application même quand le fichier est illisible ?
\item Pourquoi une forme d'onde s'affiche-t-elle en min/max et non en moyenne ?
\item Pourquoi \texttt{lo} est-il initialisé à $1$ et \texttt{hi} à $-1$ ?
\item Combien de sommets pour un rectangle plein, et pourquoi pas quatre ?
\item Pourquoi la cadence vidéo emploie-t-elle un \texttt{while} et non un
      \texttt{if} ?
\item Où se branche un traitement d'image, et pourquoi pas ailleurs ?
\end{enumerate}
\end{nkquestions}

\begin{nkpratique}
Construisez votre lecteur audio en suivant les étapes 1 à 3 et 9, sur
\texttt{NkCanvasGuiApp}.

Il doit : accepter un fichier en argument, afficher un message lisible quand il
n'y arrive pas, dessiner la forme d'onde en min/max, jouer le son, et montrer un
curseur de lecture qui avance.

\textbf{Aucun appel à \texttt{NkWindow} ni à \texttt{NkRenderWindow} dans votre
code.} Comptez les lignes, et comparez aux 263 du lecteur du dépôt : dites
lesquelles ont disparu et pourquoi.

Puis ajoutez la vidéo : étapes 4 à 6.
\end{nkpratique}

\begin{nkdemo}
Prouvez que la moyenne ne convient pas.

Ajoutez un second mode d'affichage qui dessine la \emph{moyenne} de chaque
tranche au lieu du min/max, et une touche pour basculer. Chargez un morceau réel
et capturez les deux images.

\textbf{Ce que la démonstration doit établir} : non pas que la seconde est
« moins jolie », mais qu'elle est \emph{plate} --- donc qu'elle ne distingue pas
un passage fort d'un passage faible. Mesurez-le : affichez l'écart entre la plus
grande et la plus petite valeur affichée, dans les deux modes.
\end{nkdemo}

\begin{nklogique}
Une vidéo à 25 images par seconde dure 4 minutes. Le décodage d'une image prend
en moyenne 12 ms, le téléversement 2 ms, le dessin 1 ms.

\begin{enumerate}
\item La lecture tiendra-t-elle la cadence ? Justifiez par le calcul.
\item Une image sur cinquante prend 300 ms à décoder (image-clé). La lecture
      saccade-t-elle ? Que faudrait-il changer, \emph{sans} accélérer le
      décodeur ?
\item Vous ajoutez un filtre couleur qui coûte 9 ms par image. Où le mettre ---
      et si vous n'aviez le droit de l'appliquer qu'une image sur deux, la vidéo
      serait-elle acceptable ? Argumentez sur ce que voit l'\oe il, pas sur les
      chiffres seuls.
\end{enumerate}
\end{nklogique}
